% ==============================================================================
% sections/alternatives.tex
% Created: 2026-05-08
% ==============================================================================

\subsection{AR(2): second-order Markov dynamics}\label{subsec:alt:ar2}

The baseline model assumes that the latent employment process is first-order
Markov. If employment exhibits genuine duration dependence---if the probability
of finding a job declines with the length of a non-employment spell, or if
the probability of losing a job declines with seniority---then the first-order
Markov restriction could be misspecified, and the EM estimator might attribute
duration-dependent dynamics to misclassification.

Tables~\ref{tab:ar2_implied} and~\ref{tab:ar2_params} report the results of
the AR(2) extension, which uses four consecutive waves and allows the transition
probabilities to depend on both the current and lagged true employment state.
The four-wave linked sample contains approximately 317,000 observations,
somewhat smaller than the three-wave sample because a fourth valid wave is
required.

The AR(2) estimates reveal substantial and statistically significant duration
dependence. Among non-employed workers, those who were also non-employed in the
prior quarter have a true entry probability of only 4.33\%, compared with
30.61\% for those who were employed in the prior quarter. This large difference
($\hat{\theta}_{01} - \hat{\theta}_{00} = 26.28$ percentage points) indicates
strong negative duration dependence in non-employment: workers who have been
non-employed for at least two consecutive quarters are much less likely to find
a job in the current quarter than those who recently entered non-employment.
Among employed workers, those who were also employed in the prior quarter have
a separation probability of only 4.19\%, compared with 24.57\% for those who
were non-employed in the prior quarter, indicating that recent entrants into
employment are much more likely to exit in the subsequent quarter.

Despite this rich duration dependence, the estimated misclassification probability
in the AR(2) model is $\hat{\pi} = 1.14\%$---positive, statistically significant,
and large enough to imply meaningful correction to the observed transition
matrix. Although this estimate is lower than the baseline 2.98\%, the
persistence of a positive misclassification probability confirms that duration
dependence alone cannot explain the Markov violations in the data. The two
mechanisms---duration dependence and misclassification---are jointly identified
and coexist in the AR(2) specification.

The AR(2) results also imply that the strong negative duration dependence in
the baseline model may be partly an artefact of the first-order Markov
restriction: once longer-order dynamics are accommodated, the apparent
duration-dependence-from-misclassification confound is resolved, and both
mechanisms have distinct empirical signatures.

\subsection{Tenure contamination model}\label{subsec:alt:tenure}

Tables~\ref{tab:eps_implied}, \ref{tab:eps_params}, and~\ref{tab:eps_tenure}
report the results of the tenure contamination model. This specification
exploits the joint distribution of employment status and job tenure to provide
an independent structural identification channel for misclassification.

The estimated contamination probability is $\hat{\varepsilon} = 24.23\%$: nearly
one in four tenure observations in the sample is estimated to be
``contaminated,'' in the sense that it is drawn from the wrong employment state's
tenure distribution. This high contamination rate reflects both genuine tenure
recall error and the structural implications of employment state misclassification
for flanking tenure observations. Given this contamination rate, the estimated
misclassification probability is $\hat{\pi} = 8.64\%$ under the static
specification and somewhat higher under the CTMC-linked specification.

The tenure contamination model recovers mean employment spell lengths from the
estimated transition rates. Table~\ref{tab:eps_tenure} reports that the implied
mean employment spell length is approximately 5.85 years (70.2 months), with
$P(\text{clock-consistent}) = 75.77\%$---the share of employment-tenure
sequences in the sample that are consistent with normal tenure clock accumulation
(i.e., where tenure increases by approximately one quarter per wave). This
proportion substantially exceeds what would be expected if every employment
observation were drawn independently; the high clock-consistency rate validates
the exponential duration model as a reasonable approximation.

We note that the misclassification estimate from the tenure contamination model
is higher than the baseline. This is partly a consequence of the longer
observation window (tenure flanks span multiple waves) and partly because the
exponential duration assumption introduces additional functional form restrictions.
We interpret the tenure contamination results as corroborating the direction and
approximate magnitude of the baseline finding---misclassification is present,
economically substantial, and structurally identified---rather than as a precise
alternative estimate of $\pi$.

\subsection{Matching algorithm sensitivity}\label{subsec:alt:matching}

Table~\ref{table_matching_implied} reports sensitivity of the results to the
strictness of the wave-to-wave matching algorithm. The baseline matching
algorithm links individuals using household and person identifiers only.
The ``stricter'' matching algorithm additionally requires agreement on gender,
race, and a time-invariant household characteristic; the ``strictest'' algorithm
imposes further agreement on an additional set of stable characteristics.
Stricter matching reduces the sample size by discarding potential false matches
(i.e., observations that look like the same individual across waves but may be
different individuals who were incorrectly linked by the identifier algorithm).

As matching becomes stricter, both the raw (no misclassification) transition
rates and the EM-corrected transition rates decline, and the estimated
misclassification probability falls from $\hat{\pi} = 2.82\%$ under the
baseline matching to $\hat{\pi} = 2.25\%$ under the strictest matching. This
pattern is consistent with the interpretation that some of the Markov violations
attributed to misclassification under the baseline matching reflect genuine
mismatches: when a non-employed person is incorrectly linked to an employed
observation from another household member, the resulting spurious transition
inflates measured entry and exit rates and creates apparent Markov violations
that the EM algorithm would otherwise attribute to misclassification.

Crucially, however, even the strictest matching algorithm retains a substantial
misclassification probability of 2.25\%. Matching errors are neither the
primary nor the sole source of the Markov violations in the data; genuine survey
measurement error in employment status contributes independently. This design-based
robustness check, which varies identification assumptions (how the panel is
constructed) rather than the structural model, strongly supports the main finding.

The implied entry and exit rates under the strictest matching and symmetric
misclassification correction (2.36\% and 2.25\% respectively) are slightly lower
than the baseline 2.93\% and 3.00\%, but the overall conclusion is unchanged:
true employment transition rates in South Africa are approximately three times
lower than the raw transition matrix would suggest, and employment is considerably
more persistent than previously characterised.
